% from ega0-19-fr.tex
\oldpage[0\textsubscript{IV-1}]{87}

tout $B$-module topologique discret $L$, annulé par un idéal ouvert de $B$, on
ait (cf. (18.5.1))
\[
  \operatorname{Exalcotop}_A(B,L)=0.
\]
\end{proposition}

\begin{proof}
Soit $(\mathfrak R_\lambda)$ un système fondamental décroissant d'idéaux
ouverts de $B$ et posons $B_\lambda=B/\mathfrak R_\lambda$.
Considérons une $A$-algèbre topologique discrète $C$ et un idéal
$\mathfrak J$ de $C$ tel que $\mathfrak J^2=0$, de sorte que $C$ est une
$A$-extension de $C/\mathfrak J$ par $\mathfrak J$~; supposons donné un
$A$-homomorphisme $u:B_\lambda\to C/\mathfrak J$ et formons la
$A$-extension $E_\lambda$ de $B_\lambda$ par $\mathfrak J$, image
réciproque de $C$ par l'homomorphisme $u$ (18.2.5)~; c'est une $A$-algèbre
topologique pour la topologie discrète. Si
$\operatorname{Exalcotop}_A(B,L)=0$, il existe par définition (18.5.1) un
$\mu$ tel que $\mathfrak R_\mu\subset\mathfrak R_\lambda$ et que
l'extension image réciproque $E_\mu$ soit $A$-triviale~; mais cela signifie
(18.1.6) qu'il existe un $A$-homomorphisme continu
$v:B_\mu\to E_\lambda$ tel que l'homomorphisme canonique
$B_\mu\to B_\lambda$ se factorise en
$B_\mu\xrightarrow{v}E_\lambda\to B_\lambda$~; on en conclut
aussitôt que $B_\mu\to B_\lambda\xrightarrow{u}C/\mathfrak J$ se
factorise en
$B_\mu\xrightarrow{v}E_\lambda\to C\to C/\mathfrak J$, et cela
prouve, en vertu de (19.3.1) et (19.4.3), que $B$ est une $A$-algèbre
formellement lisse. La réciproque est immédiate, en appliquant (19.3.1) au cas
où $C$ est une $A$-algèbre topologique qui est une $A$-extension de
$B_\lambda$ par $L$, et à l'homomorphisme identique
$B_\lambda\to B_\lambda=C/L$.
\end{proof}

Lorsque $A$ et $B$ sont des anneaux discrets, le critère de lissité formelle
(19.4.4) se réduit à
\begin{equation}
\tag{19.4.4.1}
\label{0.19.4.4.1-fr}
\operatorname{Exalcom}_A(B,L)=0\qquad\text{pour tout $B$-module $L$};
\end{equation}
autrement dit, \emph{toute $A$-extension commutative de $B$ par un
$B$-module doit être $A$-triviale}.

\begin{corollary}[19.4.5]
\label{0.19.4.5-fr}
Soient $A$ un anneau topologique, $B$ une $A$-algèbre topologique
(commutative).
\begin{enumerate}
  \item[(i)] Supposons que $B$ soit une $A$-algèbre formellement lisse.
  Alors, pour tout idéal ouvert $\mathfrak R$ de $B$, tout
  $(B/\mathfrak R)$-module $L$ et tout $2$-cocycle de Hochschild
  $A$-bilinéaire et symétrique $f$ de
  $(B/\mathfrak R)\times(B/\mathfrak R)$ dans $L$ (18.4.3), il existe un
  idéal ouvert $\mathfrak R'\subset\mathfrak R$ tel que, si
  $\varphi:B/\mathfrak R'\to B/\mathfrak R$ est l'homomorphisme
  canonique, $f\circ(\varphi\times\varphi)$ soit un $2$-cobord de
  Hochschild $A$-bilinéaire de
  $(B/\mathfrak R')\times(B/\mathfrak R')$ dans $L$.
  \item[(ii)] Si $B$ est un $A$-module formellement projectif (19.2.1) et si
  la condition de (i) est vérifiée, $B$ est une $A$-algèbre formellement
  lisse.
\end{enumerate}
\end{corollary}

\begin{proof}
\begin{enumerate}
  \item[(i)] Le $2$-cocycle $f$ définit une $A$-extension de Hochschild $C$
  de $B/\mathfrak R$ par $L$ (18.4.3). Appliquant (19.3.1) à $C$, à l'idéal
  de carré nul $L$ de $C$ et à l'homomorphisme identique
  $B/\mathfrak R\to B/\mathfrak R=C/L$, on en déduit la condition (i) en
  vertu de (18.4.3).
  \item[(ii)] Appliquons le critère (19.4.3), en considérant un idéal ouvert
  $\mathfrak R$ de $B$, un idéal ouvert $\mathfrak J$ de $A$ tel que
  $\mathfrak JB\subset\mathfrak R$, et une $(A/\mathfrak J)$-extension $E$
  de $B/\mathfrak R$ par $L$. Comme $B$ est un $A$-module formellement
  projectif, l'application $A$-linéaire continue canonique
  $\rho:B\to B/\mathfrak R$ se factorise en
  $B\xrightarrow{w}E\to B/\mathfrak R$, où $w$ est une application
  $A$-linéaire continue (19.2.1), qui se factorise donc elle-même en
  $B\to B/\mathfrak R'\xrightarrow{w'}E$ où $\mathfrak R'$ est un idéal
  ouvert de $B$ contenu dans $\mathfrak R$~; remplaçant $\mathfrak J$ par
  un idéal plus petit, on peut supposer que
  $\mathfrak JB\subset\mathfrak R'$. Alors l'image réciproque par
  l'homomorphisme canonique $B/\mathfrak R'\to B/\mathfrak R$ de
  l'extension $E$ de $B/\mathfrak R$ par $L$, est une
  $(A/\mathfrak J)$-extension \emph{de Hochschild} $E'$ de
  $B/\mathfrak R'$ par $L$~;

  